% GENERATED from paper/draft/06_results.md by build.py — do not edit
\section{Results}

Three decision classes are tested (§5.1). We call them the \emph{discrepancy class} (c3 in the tables: triage of a
disagreement between the beneficial-ownership register and other evidence), the \emph{freshness class} (c5: whether a
customer file must be reopened, refreshed or left alone) and the \emph{onboarding class} (c6: an assembled onboarding
file that the agent may only escalate or complete). A \emph{configuration} is a model with its protocol (§5.3). The
\emph{common core} is the set of 60 calibration, 30 held-out and 60 shifted base episodes per class that every
configuration shares (amendment A4); the analysis as registered uses all data, and we report both. Results are scored
under three answer keys: the \emph{registered key} (the oracles' answers as pre-registered), the registered key
\emph{without the two contested arms} (amendment A1) and the \emph{reviewed key} (the registered key with the author's
verdicts applied, A13).

\subsection{Delegation on interventions versus track records (the delegation test, H3)}

\textbf{The registered test favours DSA.} For each configuration and coverage level \(k \in \{1, 2\}\), each rule
delegates the \(k\) classes with the lowest certified upper bound, and we compare the outcome error on the shifted
episodes of the delegated classes (§5.6). With 14 configurations there are 28 such comparisons; a comparison is
informative only when the two rules delegate different classes. In the registered analysis on all data, the rules differ
in 10 comparisons. DSA has the lower shifted error in 9 and the higher in none (one comparison is equal), and the
mixed-effects model gives \(\beta_{\mathrm{DSA}} = -0.36\) (standard error, SE, 0.11; odds ratio 0.70; one-sided
\(p < 0.001\)), where \(\beta_{\mathrm{DSA}}\) is the model's coefficient for the DSA rule and a negative value favours
DSA. The common core, reported with equal prominence, agrees: DSA is lower in 9 of 11 differing comparisons and higher
in none (sign test \(p = 0.002\)), with \(\beta_{\mathrm{DSA}} = -0.38\) (SE 0.11; \(p < 0.001\)). Because the model
treats a shifted episode that enters both rules and both values of \(k\) as independent observations, we also resampled
configurations with replacement (400 resamples): the bootstrap standard errors are 0.12 on all data and 0.14 in the
common core, the one-sided \(p\)-values 0.002 and 0.003, and no resample gives \(\beta_{\mathrm{DSA}} \ge 0\). By the
registered criterion, the delegation test passes; §8 reports the result under all 60 logged readings of the policy.

\textbf{The advantage rests on one contested arm.} A policy audit carried out during the pilot, after interim inspection
of the first calibration data, had flagged two intervention arms whose correct answers are contested (amendment A1, post
hoc): the removal of the register excerpt in the freshness class, where the policy does not say that a file must contain
one, and an inverted notarial deed in the discrepancy class. Excluding both leaves 7 differing comparisons, 4 : 3, in
both data sets (\(\beta_{\mathrm{DSA}} = -0.09\), \(p = 0.19\) in the common core; \(-0.08\), \(p = 0.20\) on all data).
The first arm carries the result: dropping the removal arm in every class has the same effect (common core, \(p = 0.38\)
in the leave-one-out check), dropping the inversion arm does not (\(p = 0.004\)), and in the two reviewed cases of the
inverted deed the author's blind answers equalled the oracles' (A13). Rescoring with the reviewed key, described next,
gives 3 : 1 in 8 differing comparisons on all data (\(\beta_{\mathrm{DSA}} = -0.07\), \(p = 0.28\)) and 3 : 1 in 9 in
the common core (\(-0.08\), \(p = 0.28\)).

\begin{table*}
\caption{The delegation test under the registered key, without the contested arms, and under the reviewed key.
``Comparisons'': (configuration, \(k\)) comparisons in which DSA has the lower : higher shifted error, out of those in
which the two rules differ. \(\beta_{\mathrm{DSA}}\) with, in parentheses, the one-sided \(p\) from the mixed-effects
model and the one-sided \(p\) from a bootstrap over configurations (400 resamples;
\texttt{paper/pilot/bootstrap\_h3.md}).}
\small
\begin{tabular}{@{}
  >{\raggedright\arraybackslash}p{(\textwidth - 8\tabcolsep) * \real{0.2542}}
  >{\raggedright\arraybackslash}p{(\textwidth - 8\tabcolsep) * \real{0.1314}}
  >{\raggedright\arraybackslash}p{(\textwidth - 8\tabcolsep) * \real{0.2415}}
  >{\raggedright\arraybackslash}p{(\textwidth - 8\tabcolsep) * \real{0.1314}}
  >{\raggedright\arraybackslash}p{(\textwidth - 8\tabcolsep) * \real{0.2415}}@{}}
\toprule\noalign{}
\begin{minipage}[b]{\linewidth}\raggedright
Analysis
\end{minipage} & \begin{minipage}[b]{\linewidth}\raggedright
All data: comparisons
\end{minipage} & \begin{minipage}[b]{\linewidth}\raggedright
All data: \(\beta_{\mathrm{DSA}}\) (\(p\))
\end{minipage} & \begin{minipage}[b]{\linewidth}\raggedright
Common core: comparisons
\end{minipage} & \begin{minipage}[b]{\linewidth}\raggedright
Common core: \(\beta_{\mathrm{DSA}}\) (\(p\))
\end{minipage} \\
\midrule\noalign{}

Registered key & 9 : 0 of 10 & −0.36 (\textless{} 0.001; 0.002) & 9 : 0 of 11 & −0.38 (\textless{} 0.001; 0.003) \\
Contested arms excluded (A1) & 4 : 3 of 7 & −0.08 (0.20; 0.26) & 4 : 3 of 7 & −0.09 (0.19; 0.27) \\
Reviewed key (A13) & 3 : 1 of 8 & −0.07 (0.28; 0.16) & 3 : 1 of 9 & −0.08 (0.28; 0.17) \\
\bottomrule
\end{tabular}
\end{table*}

\textbf{The author's review of the key.} Late in the pilot the author reviewed 21 of 30 sampled cases, drawn where
models disagreed with the oracles, so the result is not an error rate of the oracles (A13). He answered them in three
steps: 11 blind, on paper; 9 blind after the four open readings had been discussed; and one after the second batch's
answers had been shown to him. His answers equalled the oracles' admissible sets in 14 of the 21 (6 of 11 in the first
batch; 7 of 9 in the second, where the two differences were a date copied wrongly onto his own board and a case in which
he listed three actions where the oracles admit one; and the last case), and he reports that elements such as the
decision date came into view only after several cases, so later answers are not independent of earlier ones. He first
judged the oracles wrong in one case: exception R2(b) requires a source independent of the register's controlling
authority, the documents name issuers but never their controlling authorities, and both oracles decide independence from
a field the agents do not see, so he read the exception as never established. He withdrew this verdict the next day,
after answering blind a case of the same structure in the oracles' sense, on the ground that a register and a court are
two independent bodies even though no document says who controls them. The reviewed key therefore keeps the oracles'
reading of the exception; the legal text that R2(b) paraphrases has no independence condition at all (§8). He judged
both readings defensible in three cases, which turn on two questions the policy leaves open: whether filing annual
accounts is a ``change'' that requires reopening a file (R3), and whether a file must contain a register excerpt at all
(R3, R5). The two oracles were written from one specification and share these readings, so their agreement did not test
them. Together the two readings change the correct answer for 44 task variants in the common core. The clearest case is
the freshness class. In the common core the track record certified it for GPT-6 Luna, GPT-6.1 Terra, GPT-6.1 Sol and
Qwen 3.8 27B, and DSA did not. Under the registered key these configurations then erred on 17--27\% of the shifted
freshness episodes, which made the track record look unsafe. Under the reviewed key they erred on 0--2\%: almost every
one of these errors was reopening a file after an annual-accounts filing that the shifted split had added as a decoy.

\textbf{What holds under every key.} Under each key, the LLM judge passes most of the correct-outcome decisions that the
interventions mark as resting on the wrong basis: 63--70\% on all data and 59--68\% in the common core, pooled over the
configurations it assessed (the judge test; §6.2). It also flags 17--18\% of the decisions that pass every intervention,
so the two methods disagree in both directions, and the pilot has no third standard to say which is right. As delegation
rules, the two judge-based rules had a higher shifted error than DSA in most comparisons in which they differed, under
every key. In about two thirds of these comparisons, however, DSA delegated the onboarding class and the judge-based
rule did not, and in that class the judge mostly reads the ``current register excerpt'' of R4 as an excerpt of the
beneficial-ownership register, unlike the oracles and the agents (interim observation 2 in the registration log). These
comparisons are descriptive. In the common core, the readiness gate, a two-arm rule defined for the API configurations,
had the lower shifted error in 2 of 6 differing comparisons under the registered key (the other 4 equal) and split 1 : 1
under the reviewed key (§6.2).

\textbf{What each rule certified.} By the registered criterion the delegation test passes; §7.1 gives our reading of it.
Under the registered key the DSA rule certified only the onboarding class, in 7 of 14 configurations, at 60 episodes per
class; in that class escalation is always an admissible answer. In the common core the track record certified 16
classes, with a shifted error of 5.7\% in them. Without the contested arms DSA certified 11 classes with a shifted error
of 5.8\%, about the same as the track record's; under the reviewed key DSA certified 10 classes with one shifted error
in 600 episodes and the track record 16 with three in 960 (\texttt{paper/pilot/construct\_validity.md}, which also gives
the all-data figures). Eight of the eleven differing comparisons in the common core involve a tie at the boundary of the
delegated set, resolved by the expected-error rule of A7. Restricted to the discrepancy and freshness classes, in which
delegation would mean autonomy over a decision, the comparison at \(k = 1\) gives 6 : 0, 2 : 2 and 1 : 1 differing
comparisons under the three keys (with two further equal comparisons under the reviewed key). Two further caveats apply.
The reviewed key rests on one reviewer, who saw exploratory results before giving three of his four verdicts and who
withdrew one of them after a later case, knowing the direction of the change (A13). The configuration bootstrap of Table
4 (first proposed at the A12 audit, when an interim run on 11 configurations had given a standard error of 0.149 against
the model's 0.126) leaves every reading unchanged; the sign test's \(p\)-values are descriptive (§5.7).

\subsection{Reading the trace versus intervening on the evidence (the judge test, H1′, and the declared-grounds test,
H2)}

\textbf{The judge passes most wrong-basis decisions.} The judge test asks what share of the correct-outcome decisions
that fail an interventional basis test a procedure audit passes. The judge (Claude Haiku 5.5, rubric v2 after one
revision, amendment A3) read the factual decision, its explanation and its declared grounds for every agent with tools
and for the one-pass model; decision models return no trace to read. Pooled over the eleven judged configurations, the
audit passed 322 of the 474 correct decisions that fail a basis test in the common core (67.9\%) under the registered
key, 169 of 288 (58.7\%) without the contested arms and 205 of 337 (60.8\%) under the reviewed key; on all data, 70.3\%,
63.3\% and 64.7\%. It also flagged 17--18\% of the correct decisions that pass every intervention: 28.5\% in the
onboarding class against 6--12\% in the other two, because there it reads the ``current register excerpt'' of R4 as an
excerpt of the beneficial-ownership register (§6.1). Per configuration, under the registered key, the judge passed
74--100\% of the basis failures of the six API agents (all 25 for GPT-6.1 Sol, 18 of 19 for Claude Opus 5.5) and
37--56\% of those of three local agents with tools (Gemma 4 26B, gpt-oss 120B, Mistral Small 3.2) and the one-pass
model, whose explanations are shorter and whose failures are more often outright non-reactions (§6.3); the fourth local
agent, Qwen 3.8 27B, which rarely fails by non-reaction, sits with the API agents at 85\%.

\textbf{A stronger judge does no better.} The main alternative explanation, that Haiku is weaker than several of the
models it judged, was tested before any verdict of the second judge was produced (amendment A10). GPT-6.1 Sol judged 199
correct-outcome decisions of the Claude and GPT-6 agents, drawn in strata with and without an interventional basis
failure, with the same rubric. Its Youden's \(J\) (sensitivity minus false-positive rate) was −0.07 against +0.07 for
Haiku, difference −0.14 with a 95\% bootstrap interval of {[}−0.24, −0.04{]}; it flagged none of the 40 basis failures
of the OpenAI agents, and its mean grounds score was 4.88 for decisions with a basis failure against 4.77 without. The
two judges agreed on 87\% of verdicts (Cohen's \(\kappa = 0.31\)). By the reading fixed in advance (\(J < 0.15\)), the
``convincing explanations'' account is supported: the explanations of decisions that do not depend on the decisive
evidence read as well as those that do, to two judges from two developers. Relabelling the same 199 decisions under the
other two keys changes little: without the contested arms 38 of them count as basis failures instead of 99, and Sol's
\(J\) is −0.03 against +0.10 for Haiku; under the reviewed key 52 count, and Sol's \(J\) is −0.06 against +0.07; the
difference is −0.13 in both cases, with bootstrap intervals {[}−0.25, −0.03{]} and {[}−0.24, −0.03{]}
(\texttt{runs/judge\_check/result\_keys.txt}).

\textbf{Declared grounds carry a signal, but a weak one.} The declared-grounds test compares the interventions with a
static check of what the agent itself named as decisive: the judge's verdict on whether the documents the agent declared
as decisive are those that settle the episode under the policy. This check flagged 65\% of the interventional basis
failures and 35\% of the decisions that pass every intervention (common core, registered key; 56\% and 41\% without
contested arms; 59\% and 39\% under the reviewed key). Taken as a detector of what the interventions find, it separates
better than the judge's basis verdict (sensitivity 0.65 against 0.32, false-positive rate 0.35 against 0.17), but it
flags a third of the decisions whose basis the interventions confirm, mostly for naming too many or too few documents. A
mechanical version of the check, which compares the declared decisive documents with the document the generator defines
as decisive, does not separate at all: any difference between the two sets flags 62\% of the basis failures and 84\% of
the decisions that pass every intervention, and a decisive document left undeclared flags 13\% and 26\%
(\texttt{paper/pilot/declared\_grounds\_check.md}). Agents name the register entry together with the evidence that
confirms or contradicts it, and the generator's list names the settling document alone, so what separates is the judge's
reading of the declaration against the policy, not a match with a fixed list. The test is met in the weak sense that
interventions and declared grounds disagree in both directions and the interventions find failures that the declared
grounds do not show.

\textbf{As delegation rules.} The judged-process rule certified five classes in the common core under the registered key
and the same five under the reviewed key (the discrepancy class for Opus and Sol, the freshness class for Luna, Sol and
Terra); the judged-grounds rule, with its threshold chosen on calibration (3.25 on a five-point scale), granted every
class to every API agent and to Gemma 4 26B and Qwen 3.8 27B, including the onboarding class for GPT-6.1 Terra with a
calibration outcome error of 46.7\%, and withheld every class but one from the other local models: mean grounds scores
separate verbose from terse configurations, not reliable from unreliable classes (interim observation 3). At matched
coverage, DSA had a lower shifted error than the judged-process rule in 11 of 13 differing comparisons under the
registered key (none higher, two equal), 7 of 10 (one higher) without contested arms and 6 of 11 (none higher, five
equal) under the reviewed key, and than the judged-grounds rule in 13 of 14, 9 of 11 (one higher) and 8 of 12 (none
higher); as §6.1 notes, about two thirds of these wins are comparisons in which DSA delegates the onboarding class and
the judge-based rule does not. The readiness gate, defined only for the API configurations, had the lower shifted error
in 2 of 6 differing comparisons under the registered key, with the other 4 equal, and split 1 : 1 of 6 under the
reviewed key; the pilot gives no evidence that the full intervention family delegates better than this two-arm rule. Two
further baselines could not certify at this sample size, which the design implies rather than the data: a certified
confidence threshold needs about 75 error-free decisions above the threshold at \(n = 60\), and a simulated reviewer who
errs with probability 0.05 grants even a perfect agent with probability at most 0.19 (A12). The decision models'
confidence separates correct from incorrect outcomes (AUROC 0.55--0.78) about as well as it separates right from wrong
basis among correct outcomes (0.56--0.84 under the registered key, 0.56--0.81 under the reviewed key), with Basal
2.0-max the highest on the basis measure.

\subsection{Construct checks (CV1, CV2) and the known-groups check}

\textbf{In this pilot DSA is largely, not wholly, reducible to outcome accuracy (CV1).} Across the 42 (configuration,
class) cells of the common core, the share of correct outcomes that fail a basis test rises with the outcome error:
Pearson \(r = 0.88\) (Spearman \(\rho = 0.79\)) under the registered key, 0.94 (0.94) without the contested arms and
0.93 (0.91) under the reviewed key; a linear regression on \(\ell_{\text{out}}\) leaves 23\%, 11\% and 14\% of the
variance unexplained, and 17--22\% across configurations with classes pooled
(\texttt{paper/pilot/construct\_validity.md}). The pre-registered expectation that DSA is not reducible to outcome
accuracy is therefore only weakly met: configurations that err more also fail more basis tests when they are right, and
the independent part is between a tenth and a quarter of the variance. The residual sits where the construct predicts.
Agents with no outcome error in a class still fail basis tests in 8--17\% of discrepancy episodes (Claude Haiku, Opus
and Sonnet 5.5, GPT-6.1 Sol) and in 33--40\% of freshness episodes (GPT-6 Luna, GPT-6.1 Terra and Sol, Qwen 3.8 27B;
Sonnet, with an outcome error of 3\%, in 43\%), and for these agents a large share of the basis failures, for Sonnet,
Opus, Sol, Luna and Qwen most of them, lie on the contested arms (26 of 35 for Sonnet, 14 of 19 for Opus, 20 of 25 for
Sol, 24 of 44 for Luna, 22 of 46 for Terra, 14 of 35 for Haiku and 27 of 39 for Qwen 3.8 27B), whereas for the local and
decision models few do (7 of 70 for gpt-oss 120B, 6 of 53 for Basal) and the failures are mostly non-reactions: Liquid
d1 kept its factual action in 100\% of the discrepancy-class interventions that made it inadmissible, Basal 2.0-max in
81\% and Mistral Small 3.2 in 51\%, against 0--7\% for the API agents. Whether the frontier models' residual is
dependence on evidence or disagreement with the key is what §6.1 leaves open.

\textbf{Where the variance lies (CV2).} Decomposing the unit-level basis loss of the 2,520 core calibration units by
sums of squares, 15\% lies between configurations, 20\% between classes, 13\% in their interaction and 52\% within cells
(registered key; 27\%, 10\%, 22\% and 42\% without contested arms; 26\%, 12\%, 20\% and 43\% under the reviewed key).
Most of the variance is between episodes of one configuration and class, which is where the certification bound operates
and where a sample of 60 is small. The seed component can be measured only for Claude Haiku 5.5, which ran three times:
its runs disagree on the basis loss in 22\% of episodes, and 29\% of the variance of its basis loss lies between runs of
the same episode (54\% without contested arms, 55\% under the reviewed key). For the thirteen configurations that ran
once, sampling noise and dependence on evidence cannot be separated (§8).

\textbf{Known groups.} The five scripted agents largely behaved as predicted in expectations written down before the run
(file \texttt{known\_groups.md}). The oracle agent is clean under every measure. The memorising agent, which replays
stored answers for episodes it has seen, has zero calibration outcome error in all three classes, so the track-record
rule delegates all of them, while DSA withholds the discrepancy and freshness classes (basis loss 80\% and 100\%) and
the agent's held-out and shifted errors in those classes are 43--60\%: a perfect track record with no dependence on
evidence, the case that motivates the paper (§1.2). The always-escalate agent is clean in the onboarding class and fails
the other two under both rules. The independence-blind agent errs only where independence decides (discrepancy class:
outcome error 10\%, basis loss 25\%, shifted error 25\%; the other classes clean). The cue-following agent flips under
placebo and sham arms in 28\% of discrepancy and 16\% of freshness arms and never in the onboarding class. One
expectation was not met: the copy arm did not reveal the independence-blind agent (copy flip 0\%), because adding a
dependent copy of the decisive document does not change the oracles' answer in these episodes; the deficit shows only
through the add-source arms and the shifted split. The check validates the code and the measures on deficits we planted;
it does not bear on the delegation test, and in the onboarding class both rules certify every planted agent, because
escalation is always admissible there.
